\documentclass[../mathNotesPreamble]{subfiles}

\providecommand{\relscalefact}{1.4}
\begin{document}
\relscale{\relscalefact}
  \section{9.1:  Introduction to Probability}
  %TODO define sample space, P(E) equally likely, number of elements in A (?)
  \begin{defn*}
    A \textbf{sample space} is the set of all possible outcomes of a random process or experiment. An \textbf{event} is a subset of a sample space.
    
    If $S$ is a finite sample space in which all outcomes are equally likely and $E$ is an even in $S$, then the \textbf{probability of $E$} denoted $P(E)$, is
      \[P(E)=\frac{\textnormal{number of outcomes in } E}{\textnormal{total number of outcomes in } S}\]
    \vspace*{0.5\baselineskip}
    
    For any finite set $A$, $N(A)$ denotes the number of elements in $A$. Thus,
      \[P(E)=\frac{N(E)}{N(S)}\]
  \end{defn*}
  %TODO deck of cards example (9.9.1), rolling dice example (9.1.2)
  \begin{ex*}
    Consider a deck of $52$ playing cards divided into four suits: 
      \raisebox{-0.6ex}{
        \hfill \scalebox{1.65}{\ensuremath{\spadesuit}} 
        \hfill \scalebox{1.65}{\ensuremath{\textcolor{red}{\varheartsuit}}} 
        \hfill \scalebox{1.65}{\ensuremath{\clubsuit}} 
        \hfill \scalebox{1.65}{\ensuremath{\textcolor{red}{\vardiamondsuit}}}
      }
     \begin{extasks}[after-item-skip=\stretch{1}](1)
       \task What is the sample space?
       \task What is the event that the chosen card is a red face card?
       \task What is the probability that the chosen card is a red face card?
     \end{extasks}
     \vspace*{\stretch{1}}
  \end{ex*}
  \pagebreak

  \begin{ex*}
    Consider rolling a pair of $6$-sided die. 
    \begin{extasks}[after-item-skip=\stretch{1}](1)
      \task Write the sample space $S$ of possible outcomes.
      \task Write the event $E$ that the numbers appearing face up have a sum of $6$ and find the probability of this event.
    \end{extasks}
    \vspace*{\stretch{1}}
  \end{ex*}
  \pagebreak

  %TODO Theorem 9.1.1 num elements in a list
  \begin{thmBox*}[The Number of Elements in a List]
    If $m$ and $n$ are integers and $m\leq n$, then there are $n-m+1$ integers from $m$ to $n$ inclusive.
  \end{thmBox*}
  %TODO example 9.1.4 (elements in subset)
  \begin{ex*}
    \mbox{}
    \begin{extasks}[after-item-skip=\stretch{1}](1)
      \task How many three-digit integers are divisible by $5$?
      \task What is the probability that a randomly chosen three-digit integer is divisible by $5$?
    \end{extasks}
    \vspace*{\stretch{1}}
  \end{ex*}
  \pagebreak
 
  %TODO example 9.1.5 (elements in array)
  \begin{ex*}
    Let $A[1], A[2],\dots, A[n]$ be a one-dimensional array, where $n$ is a positive integer.
    \begin{extasks}[after-item-skip=\stretch{1}](1)
      \task Suppose the array is cut at a middle value $A[m]$ so that two subarrays are formed:
        \[A[1], A[2],\dots, A[m] \textnormal{ and } A[m+1], A[m+2],\dots, A[n]\]
        How many elements does each subarray have?
      \task 
        What is the probability that a randomly chosen element of the array has an even subscript?
    \end{extasks}
    \vspace*{\stretch{1}}
  \end{ex*}

  \pagebreak
\end{document}